\section{Related Work}
\label{sec:related}

Atelier addresses AI-assisted research at the protocol layer, between freeform model use and autonomous-scientist systems. Autonomous-scientist platforms \citep{lu2024aiscientist, gottweis2025aicoscientist} aim to operate the research loop end-to-end, and benchmarks for ML and science agents \citep{chan2024mlebench, chen2024scienceagentbench} document that current systems remain brittle on the tasks researchers care about. Both lines illustrate the post-hoc-optimization and brittle-correctness pitfalls noted in the introduction. Atelier inverts the central design assumption. The human remains principal, the lifecycle's revision triggers route weak evidence back to method or frame revision rather than to metric selection, and manager resolution and critic packets force a separation between report and accepted evidence.

Component-level tools compose with the protocol rather than substitute for it. Literature and idea agents \citep{baek2024researchagent, skarlinski2024paperqa} produce research artifacts that Atelier can treat as worker reports requiring manager resolution. Execution frameworks \citep{yao2023react, wu2023autogen, jiang2025aide} supply patterns for how agents reason, converse, and search. Atelier uses such patterns inside bounded packets without inheriting their lack of evidence semantics. Agent-native research artifacts \citep{liu2026ara} encode finished research, while Atelier targets the upstream process and could output one as a deliverable. Across these lines, Atelier's contribution is the protocol layer that the surveyed work does not supply, while autonomous-scientist platforms substitute for the human and answer a different question.
